\scanpage{40}
\noindent\makebox[0pt][r]{\rule[-1.4ex]{.4pt}{3.7ex}\hspace{.5em}}\uline{Proposition} Soi\deleted{\unsure{ent|t}}t $(C,\sigma,<)$ une \unsure{maquette|moquette} ordonnée. Alors pour tt $x\in C$, les \deleted{\unsure{ens.|sous}} ens.
\[
\begin{aligned}
 C_{\geq x}&=\{y\in C\mid y\geq x\}\\
 C_{\leq x}&=\{y\in C\mid y\leq x\}
\end{aligned}
\]
\makebox[0pt][r]{\rule[-.4ex]{.4pt}{2.4ex}\hspace{.5em}}sont \uline{tot\textsuperscript{\unsure{t|l}}} ordonnés.

Il suffit de \inserted{le} prouver \unsure{\deleted{le}|\deleted{la}|le}\inserted{pour} premier \deleted{\unsure{notons|noter}} (par dualité), \deleted{Soit \unsure{donc}} \unsure{cel|cela} signifie \unsure{aussi|ceci} :
\[
\begin{aligned}
 &\text{Si}\quad x<y,\quad x<z,\quad \text{avec}\quad y\ne z\\
 &\hspace{1em}\text{alors}\quad (y,z)\ \text{comparables.}
\end{aligned}
\]

\noindent\begin{minipage}[c]{.38\linewidth}
\diagram{\begin{tikzcd}[ampersand replacement=\&,row sep={7.4mm,between origins},column sep={12.4mm,between origins},
execute at end picture={
\draw[line width=.6pt] ([xshift=-14mm,yshift=-1.1mm]Jfortyx.center) -- ([xshift=-10mm,yshift=.2mm]Jfortyx.center);
\draw[line width=.45pt] ([xshift=-3.4mm,yshift=2.3mm]Jfortyx.center) -- ([xshift=-1.7mm,yshift=2.8mm]Jfortyx.center);
\draw[->,line width=.6pt] ([xshift=4mm,yshift=4.2mm]Jfortyx.center) -- ([xshift=8.6mm,yshift=6.3mm]Jfortyx.center);
\draw[line width=.75pt] ([xshift=14mm,yshift=-.8mm]Jfortyx.center) -- ([xshift=6.3mm,yshift=1.9mm]Jfortyx.center);
\draw[->,line width=.85pt] ([xshift=1.2mm,yshift=3.1mm]Jfortyx.center) -- ([xshift=-8mm,yshift=6.2mm]Jfortyx.center);
}]
 |[inner sep=0pt,label=above left:{\(z\)}]| \phantom{\cdot} \& \cdot \& |[inner sep=0pt,label=right:{\(y\)}]| \cdot\\
 \& |[alias=Jfortyx,inner sep=0pt,label=below:{\(x\)}]| \cdot\arrow[ul,shorten <=2.9mm,shorten >=.3mm]\arrow[ur,shorten <=2mm,shorten >=.5mm] \&
\end{tikzcd}}
\end{minipage}\hfill\begin{minipage}[c]{.57\linewidth}
[On note que $y$ et $z$ sont non congrus mod $\sigma$, à cause de 28 a),]
\end{minipage}

Si $\deleted{\unsure{x}}y$ et $\deleted{\unsure{y}}z$ étaient incomparables; on aurait que $\sigma y$ et $z$ sont comparables, On \unsure{peu|pas|ne} peut avoir $z\leq\sigma y$ car on. aurait $x\leq\sigma y$, ce qui contredit $x<y$ par (28) a). On \deleted{\unsure{ne peut avoir}} \inserted{aurait donc} $\sigma y<z$ d'où $\sigma z<\backslash$

\clearpage
% archive sheet number: 39
